% GENERATED FILE. Edit canonical lessons, then run npm run generate:books.
% canonical-source-hash: fnv1a64:b836cf1a312413e9
% canonical-lessons: TE-C290-pagalu, TE-C290-puta, TE-C290-dasabdam, TE-C290-satabdam, TE-C290-balyam

\chapter{Daytime, Half of the day, Decade, Century, Childhood}
\label{ch:te-daytime-half-of-the-day-decade-century-childhood}
\hlchaptermodality{290}

\begin{chapteropening}
\textbf{By the end of this chapter:} \emph{I can say daytime, a half of the day, a decade, a century and childhood.}

Complete the last lesson of chapter 290: I can say daytime, a half of the day, a decade, a century and childhood.
\end{chapteropening}

\section[\texorpdfstring{\te{పగలు} (pagalu)}{పగలు (pagalu)}]{\te{పగలు} (pagalu) — daytime}

\label{lesson:TE-C290-pagalu}



\begin{quote}
Before the new one: say the Telugu for a celebration, then the Telugu for sunrise.
\end{quote}

\subsection*{You'll want to know: \te{పగలు}}
\textbf{\te{పగలు}} — \emph{pagalu} — \textquotedblleft{}daytime\textquotedblright{}.

\begin{grammarlens}[title={What you've built}]
One more word for this chapter.
\end{grammarlens}

\subsection*{Your turn}
\begin{itemize}
\raggedright
  \item \emph{Say it:} \emph{pagalu}
  \item \emph{Say it:} \emph{pagalu}, once more
  \item \emph{Say it:} say \emph{pagalu}
  \item \emph{Recall:} say the Telugu for a celebration, then the Telugu for sunrise, then say \emph{pagalu} again
\end{itemize}

\begin{tcolorbox}[breakable,colback=teal!4,colframe=teal!35!black,title={Before you move on}]
What does \te{పగలు} mean? (\textquotedblleft{}Daytime\textquotedblright{}.) Say it once more.
\end{tcolorbox}
\section[\texorpdfstring{\te{పూట} (pūṭa)}{పూట (pūṭa)}]{\te{పూట} (pūṭa) — a half of the day (a mealtime)}

\label{lesson:TE-C290-puta}



\begin{quote}
Before the new one: say the Telugu for sunrise, then the Telugu for daytime.
\end{quote}

\subsection*{You'll want to know: \te{పూట}}
\textbf{\te{పూట}} — \emph{pūṭa} — \textquotedblleft{}a half of the day (a mealtime)\textquotedblright{}.

\begin{grammarlens}[title={What you've built}]
One more word for this chapter.
\end{grammarlens}

\subsection*{Your turn}
\begin{itemize}
\raggedright
  \item \emph{Say it:} \emph{pūṭa}
  \item \emph{Say it:} \emph{pūṭa}, once more
  \item \emph{Say it:} say \emph{pūṭa}
  \item \emph{Recall:} say the Telugu for sunrise, then the Telugu for daytime, then say \emph{pūṭa} again
\end{itemize}

\begin{tcolorbox}[breakable,colback=teal!4,colframe=teal!35!black,title={Before you move on}]
What does \te{పూట} mean? (\textquotedblleft{}A half of the day (a mealtime)\textquotedblright{}.) Say it once more.
\end{tcolorbox}
\section[\texorpdfstring{\te{దశాబ్దం} (daśābdaṁ)}{దశాబ్దం (daśābdaṁ)}]{\te{దశాబ్దం} (daśābdaṁ) — a decade}

\label{lesson:TE-C290-dasabdam}



\begin{quote}
Before the new one: say the Telugu for daytime, then the Telugu for a half of the day.
\end{quote}

\subsection*{You'll want to know: \te{దశాబ్దం}}
\textbf{\te{దశాబ్దం}} — \emph{daśābdaṁ} — \textquotedblleft{}a decade\textquotedblright{}.

\begin{grammarlens}[title={What you've built}]
One more word for this chapter.
\end{grammarlens}

\subsection*{Your turn}
\begin{itemize}
\raggedright
  \item \emph{Say it:} \emph{daśābdaṁ}
  \item \emph{Say it:} \emph{daśābdaṁ}, once more
  \item \emph{Say it:} say \emph{daśābdaṁ}
  \item \emph{Recall:} say the Telugu for daytime, then the Telugu for a half of the day, then say \emph{daśābdaṁ} again
\end{itemize}

\begin{tcolorbox}[breakable,colback=teal!4,colframe=teal!35!black,title={Before you move on}]
What does \te{దశాబ్దం} mean? (\textquotedblleft{}A decade\textquotedblright{}.) Say it once more.
\end{tcolorbox}
\section[\texorpdfstring{\te{శతాబ్దం} (śatābdaṁ)}{శతాబ్దం (śatābdaṁ)}]{\te{శతాబ్దం} (śatābdaṁ) — a century}

\label{lesson:TE-C290-satabdam}



\begin{quote}
Before the new one: say the Telugu for a half of the day, then the Telugu for a decade.
\end{quote}

\subsection*{You'll want to know: \te{శతాబ్దం}}
\textbf{\te{శతాబ్దం}} — \emph{śatābdaṁ} — \textquotedblleft{}a century\textquotedblright{}.

\begin{grammarlens}[title={What you've built}]
One more word for this chapter.
\end{grammarlens}

\subsection*{Your turn}
\begin{itemize}
\raggedright
  \item \emph{Say it:} \emph{śatābdaṁ}
  \item \emph{Say it:} \emph{śatābdaṁ}, once more
  \item \emph{Say it:} say \emph{śatābdaṁ}
  \item \emph{Recall:} say the Telugu for a half of the day, then the Telugu for a decade, then say \emph{śatābdaṁ} again
\end{itemize}

\begin{tcolorbox}[breakable,colback=teal!4,colframe=teal!35!black,title={Before you move on}]
What does \te{శతాబ్దం} mean? (\textquotedblleft{}A century\textquotedblright{}.) Say it once more.
\end{tcolorbox}
\section[\texorpdfstring{\te{బాల్యం} (bālyaṁ)}{బాల్యం (bālyaṁ)}]{\te{బాల్యం} (bālyaṁ) — childhood}

\label{lesson:TE-C290-balyam}



\begin{quote}
Before the new one: say the Telugu for a decade, then the Telugu for a century.
\end{quote}

\subsection*{You'll want to know: \te{బాల్యం}}
\textbf{\te{బాల్యం}} — \emph{bālyaṁ} — \textquotedblleft{}childhood\textquotedblright{}.

\begin{grammarlens}[title={What you've built}]
One more word for this chapter.
\end{grammarlens}

\subsection*{Your turn}
\begin{itemize}
\raggedright
  \item \emph{Say it:} \emph{bālyaṁ}
  \item \emph{Say it:} \emph{bālyaṁ}, once more
  \item \emph{Say it:} say \emph{bālyaṁ}
  \item \emph{Recall:} say the Telugu for a decade, then the Telugu for a century, then say \emph{bālyaṁ} again
\end{itemize}

\begin{tcolorbox}[breakable,colback=teal!4,colframe=teal!35!black,title={Before you move on}]
What does \te{బాల్యం} mean? (\textquotedblleft{}Childhood\textquotedblright{}.) Say it once more.
\end{tcolorbox}
